\lysh{15016}{2}
\begin{alist}
\item Το ζεύγος $(0, 4)$ επαληθεύει μόνο την εξίσωση $3x + 2y = 8$ και όχι την εξίσωση $2x - y = 3$, οπότε δεν αποτελεί λύση του συστήματος.
\item Από τη δεύτερη εξίσωση του συστήματος έχουμε $2x - y = 3 \Leftrightarrow -y = -2x + 3 \Leftrightarrow y = 2x - 3$ και με αντικατάσταση στην εξίσωση $3x + 2y = 8$ έχουμε:
\[3x + 2(2x - 3) = 8 \Leftrightarrow 3x + 4x - 6 = 8 \Leftrightarrow 7x = 14 \Leftrightarrow x = 2.\]
Για $x = 2$ στην εξίσωση $2x - y = 3$ έχουμε $y = 1$.
Συνεπώς η λύση του συστήματος είναι το ζεύγος $(2, 1)$.
\item Οι συντεταγμένες του σημείου τομής των ευθειών $(\varepsilon_1), (\varepsilon_2)$ είναι η λύση του συστήματος
\[\begin{cases} 3x + 2y = 8 \\ 2x - y = 3 \end{cases},\]
δηλαδή το ζεύγος $(2, 1)$.
\end{alist}

